\section*{摘要}
\addcontentsline{toc}{section}{摘要}
本实验使用TonyPi人形机器人和TurtleBot2移动机器人，完成抱物穿行与多机跟随任务。TonyPi需要抱起物块，沿围挡通道完成左转和右转，并在终点放下；后方车辆采用不同方法保持跟随，同时处理障碍、转角和目标暂时不可见等情况。

我们大组分工完成了人形机器人路线控制以及四类跟随方案。第1组以录制动作组和视觉反馈控制抱物行走，利用固定标记识别和PnP位姿解算判断路线阶段，并研究地图定位与避障跟随；第2组使用前车背贴AR标记进行相对位姿跟随；第3组采用RGB-D目标检测、连续关联和局部地图组织实时跟随；第10组将点云跟踪首车与AR跟随后车组成三车队列，结合扫描数据调整转向。

实验中，我们共同完成场地搭建、参数调整、程序联调和资料整理。调试结果说明，动作衔接、相机视野、标签尺寸、坐标变换和速度限制都直接影响协同效果。三车队列完成了直行、转弯与顺序停车，围栏和慢行前车引起的误避障仍需要改进；部分方法也还需要补充完整场景下的测试。报告结合各组的实现、现场现象和问题分析，总结了不同跟随方法的特点与不足。

\noindent\textbf{关键词：}TonyPi；TurtleBot2；视觉跟随；动作组；地图导航；多机协同
